\section{Lecture 2: A Probabilistic View of Learning}

\subsection{Motivation}

In the previous lecture, we formulated learning as the problem of approximating an unknown function from data. However, this formulation leaves open a fundamental question:

\medskip

\noindent
\textit{Where does the data come from?}

\medskip

\noindent
In practice, data is not deterministic but arises from a stochastic process. Even if there exists an underlying functional relationship, observations are typically corrupted by noise or variability.

\medskip

\noindent
To capture this, we adopt a probabilistic framework in which both inputs and outputs are modeled as random variables.

\medskip

\noindent
\textbf{Guiding idea.}  
Learning is the problem of estimating relationships between random variables based on finite samples.

\subsection{Data as Random Variables}

We assume that data points are drawn independently from an unknown joint distribution:
\[
(x, y) \sim P,
\]
where $P$ is a probability distribution on $\mathcal{X} \times \mathcal{Y}$.

\medskip

\noindent
The dataset $\{(x_i, y_i)\}_{i=1}^n$ is thus viewed as an i.i.d.\ sample from $P$.

\medskip

\noindent
\textbf{Interpretation.}
\begin{itemize}
    \item The true relationship between $x$ and $y$ is encoded in $P$
    \item The function $f$ from the previous lecture is now understood as a conditional expectation or decision rule derived from $P$
\end{itemize}

\subsection{Risk and Expected Loss}

Given a model $f : \mathcal{X} \to \mathcal{Y}$, we define its \emph{risk} as the expected loss:
\[
R(f) = \mathbb{E}_{(x,y) \sim P} \left[ \ell(f(x), y) \right].
\]

\medskip

\noindent
\textbf{Goal.}  
Find a function $f$ that minimizes the risk:
\[
f^* = \arg\min_f R(f).
\]

\medskip

\noindent
\textbf{Bayes optimal predictor.}  
The function $f^*$ is called the \emph{Bayes predictor}.

\medskip

\noindent
\textbf{Example (squared loss).}  
If $\ell(f(x), y) = (f(x) - y)^2$, then:
\[
f^*(x) = \mathbb{E}[y \mid x].
\]

\medskip

\noindent
\textbf{Interpretation.}
\begin{itemize}
    \item The optimal prediction is the conditional expectation
    \item Learning seeks to approximate this conditional structure
\end{itemize}

\subsection{Empirical Risk Minimization}

Since the true distribution $P$ is unknown, we cannot compute $R(f)$ directly.

\medskip

\noindent
Instead, we approximate it using the empirical risk:
\[
\hat{R}_n(f) = \frac{1}{n} \sum_{i=1}^n \ell(f(x_i), y_i).
\]

\medskip

\noindent
\textbf{Empirical Risk Minimization (ERM).}
\[
\hat{f} = \arg\min_{f \in \mathcal{H}} \hat{R}_n(f),
\]
where $\mathcal{H}$ is the hypothesis space.

\medskip

\noindent
\textbf{Justification.}
\begin{itemize}
    \item By the law of large numbers, $\hat{R}_n(f)$ approximates $R(f)$
    \item Minimizing empirical risk approximates minimizing true risk
\end{itemize}

\medskip

\noindent
\textbf{Key challenge.}  
The empirical minimizer may not generalize well to unseen data.

\subsection{Conditional Distributions and Prediction}

The probabilistic framework emphasizes modeling the conditional distribution:
\[
p(y \mid x).
\]

\medskip

\noindent
\textbf{Prediction strategies.}
\begin{itemize}
    \item Regression: predict $\mathbb{E}[y \mid x]$
    \item Classification: predict $\arg\max_y p(y \mid x)$
\end{itemize}

\medskip

\noindent
\textbf{Example (classification).}  
For $y \in \{1, \dots, K\}$:
\[
f^*(x) = \arg\max_k p(y = k \mid x).
\]

\medskip

\noindent
\textbf{Insight.}  
Learning can be viewed as estimating conditional distributions or functions derived from them.

\subsection{Bayes’ Theorem}

Bayes’ theorem provides a fundamental relationship between conditional and marginal probabilities:
\[
p(y \mid x) = \frac{p(x \mid y)\, p(y)}{p(x)}.
\]

\medskip

\noindent
\textbf{Components.}
\begin{itemize}
    \item $p(y)$: prior distribution
    \item $p(x \mid y)$: likelihood
    \item $p(y \mid x)$: posterior
    \item $p(x)$: evidence
\end{itemize}

\medskip

\noindent
\textbf{Interpretation.}
\begin{itemize}
    \item Prior encodes initial beliefs
    \item Likelihood captures how data is generated
    \item Posterior updates beliefs given observations
\end{itemize}

\medskip

\noindent
\textbf{Bayesian viewpoint.}  
Learning is the process of updating beliefs about unknown quantities using observed data.

\subsection{Loss Functions and Probabilistic Modeling}

Loss functions often arise naturally from probabilistic assumptions.

\medskip

\noindent
\textbf{Negative log-likelihood.}  
If we model $p_\theta(y \mid x)$, we may define:
\[
\ell(f_\theta(x), y) = -\log p_\theta(y \mid x).
\]

\medskip

\noindent
\textbf{Result.}
\[
\min_\theta \sum_i -\log p_\theta(y_i \mid x_i)
\]
corresponds to maximum likelihood estimation.

\medskip

\noindent
\textbf{Examples.}
\begin{itemize}
    \item Gaussian noise $\rightarrow$ squared loss
    \item Bernoulli model $\rightarrow$ cross-entropy loss
\end{itemize}

\medskip

\noindent
\textbf{Insight.}  
Loss functions encode assumptions about the data-generating process.

\subsection{Uncertainty and Noise}

The probabilistic framework explicitly accounts for uncertainty.

\medskip

\noindent
\textbf{Sources of uncertainty.}
\begin{itemize}
    \item Noise in observations
    \item Incomplete information
    \item Model limitations
\end{itemize}

\medskip

\noindent
\textbf{Prediction with uncertainty.}
\begin{itemize}
    \item Instead of predicting a single value, we may predict a distribution
    \item Confidence can be quantified
\end{itemize}

\medskip

\noindent
\textbf{Insight.}  
Probability provides a natural language for reasoning under uncertainty.

\subsection{A Unifying Perspective}

The probabilistic view reframes learning as statistical estimation:

\begin{itemize}
    \item \textbf{Data:} samples from an unknown distribution
    \item \textbf{Model:} approximation of conditional structure
    \item \textbf{Objective:} minimize expected loss
\end{itemize}

\medskip

\noindent
\textbf{Core abstraction.}  
Learning is the estimation of functions or distributions from random data.

\medskip

\noindent
\textbf{Key insight.}  
The deterministic function approximation view is a special case of a broader probabilistic framework.

\subsection{Outlook}

In this lecture, we introduced a probabilistic formulation of learning, including:

\begin{itemize}
    \item data as random variables,
    \item risk and empirical risk minimization,
    \item conditional distributions,
    \item Bayes’ theorem.
\end{itemize}

\medskip

\noindent
This perspective raises further questions:

\begin{itemize}
    \item How should we choose the hypothesis space?
    \item What assumptions guide learning?
    \item How do we balance flexibility and generalization?
\end{itemize}

\medskip

\noindent
In the next lecture, we study hypothesis spaces and inductive bias, which determine the structure of the functions we learn.

\medskip

\noindent
\textbf{Guiding principle.}  
Learning is fundamentally a statistical problem: we infer structure from randomness using limited data.